% Section 9.4, from lectures/L09/appendix/c_exercises.md. The worked solutions are not
% printed; the aside says where they are.

\section{Exercises}
\label{c9:sec:exercises}

Eight, ending with the Cross-check. That one asks a design lever to keep helping, and finds that past a certain point it stops.

\begin{aside}
\textbf{How to check your work.} A \kindtag{Code} exercise is judged by the suite that ships with it: run \code{make test} at the root of the companion repository, with \code{AEL_DIR} pointing at your toolkit. A suite you have not started is reported as \code{SKIP} rather than as a failure, so the command is worth running from the first day. The hand calculations and designs are checked against the toolkit functions each one names.

\textbf{The solutions are not in this book.} They are published in full in the companion repository, in \file{lectures/L09/appendix}, including the plausible wrong answer where there is one. Read them once you have your own answers, including the ones you are unsure about, because being unsure and right is a different thing from being unsure and wrong and the solutions say which is which.
\end{aside}

% ------------------------------------------------------------------
\recall{The pair and its tail}
\label{c9:ex:1}
\begin{enumerate}
  \item What does the tail current fix, and what do the inputs decide?
  \item A pair has a 2 mA tail. What is $r_e$ on each side? Show the step that is easy to get wrong.
  \item Describe what the tail \emph{node} does for a differential input and for a common-mode input.
  \item Why does a resistance $R_{tail}$ in the tail behave as $2R_{tail}$ to one half of the pair?
\end{enumerate}

% ------------------------------------------------------------------
\recall{The two factors of two}
\label{c9:ex:2}
\begin{enumerate}
  \item The differential gain to one collector is $-R_C/2r_e$. Account for the two.
  \item There is a second factor of two available. Where, and what recovers it?
  \item Which of the two is a property of what ``differential'' means, and which is a choice?
  \item A current-mirror load raises the gain by a factor of ten in this chapter. Split that into its two independent causes and say what each is worth.
\end{enumerate}

% ------------------------------------------------------------------
\handcalc{The pair, resistively loaded}
\label{c9:ex:3}
A pair with a 2 mA tail, $R_C = 10$ kilohm, $\beta = 50$, and a 10 kilohm tail resistor.

\begin{enumerate}
  \item $r_e$, the differential gain to one collector, and the gain to both.
  \item The common-mode gain to one collector.
  \item CMRR, as a ratio and in decibels.
  \item Recompute all of parts 1 to 3 with $R_C = 100$ kilohm. Which numbers changed and which did not?
\end{enumerate}
\checkyourself{\code{differentialGain}, \code{commonModeGain}, \code{commonModeRejection}}

% ------------------------------------------------------------------
\handcalc{The large signal}
\label{c9:ex:4}
Same pair.

\begin{enumerate}
  \item The difference current at $v_d = 5$ mV, 26 mV and 100 mV.
  \item What fraction of the tail has moved at each.
  \item The input at which the small-signal answer is 1 per cent optimistic.
  \item Repeat part 3 for a 20 mA tail. Explain the result.
\end{enumerate}
\checkyourself{\code{transfer}, \code{linearRange}}

% ------------------------------------------------------------------
\design{A pair to a rejection requirement}
\label{c9:ex:5}
An instrumentation front end needs 80 dB of common-mode rejection at a 2 mA tail, with a single-ended output.

\begin{enumerate}
  \item What tail resistance does that require?
  \item What voltage would a resistor of that value drop, and what supply does that imply?
  \item State the arrangement you would actually use, and what its incremental resistance has to be.
  \item Now suppose you may take the output differentially instead. What matching would give you 80 dB with a 10 kilohm resistor tail? Comment on which of the two designs you would build.
\end{enumerate}
\checkyourself{\code{commonModeRejection}, \code{commonModeRejectionDifferential}, \code{decibels}}

% ------------------------------------------------------------------
\design{Offset}
\label{c9:ex:6}
The same pair, driven from a source of 10 kilohm on one input and 1 kilohm on the other.

\begin{enumerate}
  \item The base current in each input, and the offset voltage the \emph{imbalance} produces.
  \item The offset from a 1 per cent collector-resistor mismatch, referred to the input.
  \item The offset from 1 mV of $V_{BE}$ mismatch, referred to the input.
  \item Rank the three, then say what single change removes the largest one entirely and what it costs.
\end{enumerate}
\checkyourself{\code{inputOffset}}

% ------------------------------------------------------------------
\codeexercise{The pair}
\label{c9:ex:7}
Implement \code{ael::diffpair} to the specification in \secref{c9:sec:b7}.

\textbf{\code{commonModeRejection} must be the ratio of your own two gain functions}, not a separate formula. The chapter's central claim is that $R_C$ cancels out of that ratio, and writing it as a ratio is what makes your code demonstrate the claim rather than restate it.

\textbf{\code{linearRange} must not take the tail current as an argument.} If your derivation produced one, something cancelled that you did not cancel.

% ------------------------------------------------------------------
\crosscheck{Common-mode rejection, and the lever that reverses}
\label{c9:ex:8}
The pair of Exercise~\ref{c9:ex:3}: 2 mA tail, 10 kilohm collector resistors, 10 kilohm tail resistor to a negative rail. Find its CMRR, four ways.

\begin{enumerate}
  \item \textbf{By hand.} $(2R_{tail} + r_e)/2r_e$.
  \item \textbf{By your closed form, twice}, once with $R_C = 10$ kilohm and once with $R_C = 100$ kilohm.
  \item \textbf{By your solver.} Apply a small differential input, measure the change at one collector; apply a small common-mode input, measure the change at the same collector; divide.
  \item \textbf{Then replace the tail resistor with an ideal current source and repeat leg 3.}
\end{enumerate}

\task{What to expect}
\textbf{Legs 1, 2 and 3 agree to about a decibel}, at 52, 52 and 51. \textbf{Leg 2 gives the same number twice}, because the collector resistor cancels; if yours does not, \code{commonModeRejection} is not a ratio of your two gains.

\textbf{Leg 4 is the exercise, and it does not return infinity.} An ideal current source has infinite incremental resistance, so the expected answer is that the common-mode gain is zero. What you will get is a small, \textbf{stable, positive} common-mode gain and a CMRR of about \textbf{101 dB}.

That number is not noise: check it by changing your perturbation size by three decades and watching it stay put. \textbf{Find out where it comes from.} Two experiments settle it:

\begin{itemize}
  \item Set the Early voltage to $10^9$ and repeat. The common-mode gain should collapse to your solver's arithmetic floor.
  \item Set $\beta$ to 5000 and repeat. It should fall by about a hundred.
\end{itemize}
\textbf{Then sweep the tail.} Put the ideal source in parallel with a resistor and step that resistor from 1 megohm to 10. The closed form says CMRR rises without limit. Report what the solver says, and where the maximum is.

\textbf{Finally, answer this.} Chapter~\ref{c5:ch:transistor}'s Cross-check had two legs that agreed with each other and were both wrong about a real device. This one has a leg that is exactly right about the model and answers a question nobody asked. \textbf{What is the single statement about models that covers both?}

\textbf{If leg 3 gives a CMRR near 1}, you have applied the common-mode input to only one base.

\textbf{If leg 3's differential gain is 385 rather than 192}, you have measured the difference between the two collectors rather than one of them, which is \secref{c9:sec:b5}'s question rather than this one.
